\clearpage
\pdfbookmark[1]{Scherings Bemerkungen}{schering-remarks}
\begin{center}\textbf{Scherings Bemerkungen}\end{center}
\noindent Diese mit KI erstellte Fassung in modernem Deutsch beruht auf Ernst Scherings deutschem Original \textit{Bemerkungen} in Carl Friedrich Gauss, \textit{Werke}, Band~V, S.~637--640, herausgegeben von der Königlichen Gesellschaft der Wissenschaften zu Göttingen. Siehe das \href{https://gdz.sub.uni-goettingen.de/id/PPN236006339}{Göttinger Digitalisat} und den \href{https://archive.org/download/Werkecarlf05gausrich/page/n648}{zweiten Abdruck (1877)}.
\par\smallskip\noindent Die eigenständige englische Übersetzung \href{https://neo-classical-physics.info/uploads/3/4/3/6/34363841/schering_-_remarks_on_gauss.pdf}{\textit{Remarks}} von D. H. Delphenich beruht auf der Göttinger Ausgabe von 1867, S.~637--640.
\par\smallskip\noindent Das von Schering erwähnte Werk zur Beugung ist F. M. Schwerds \href{https://doi.org/10.3931/e-rara-16523}{\textit{Die Beugungserscheinungen aus den Fundamentalgesetzen der Undulationstheorie}}, Mannheim, Schwan \& Goetz, 1835. Das digitalisierte Exemplar liegt in der ETH-Bibliothek Zürich, Rar 2268.
\clearpage
